% --- Archon physics formalization source begin ---
% archon:physics
% archon:covers PhyXMiniProblems/problem_phyx_mini_0781.lean
% archon:source-report reports/phyx_mini/problem_phyx_mini_0781.source.json

\chapter{Physics problem phyx\_mini\_0781}
\label{ch:PhyXMiniProblems_problem_phyx_mini_0781}

\paragraph{Problem source.}
\#\# Physical scenario

A \textbackslash{}(50.0\textbackslash{}text\{ kg\}\textbackslash{}) grindstone is a solid disk \textbackslash{}(0.520\textbackslash{}text\{ m\}\textbackslash{}) in diameter. You press an ax down on the rim with a normal force of \textbackslash{}(160\textbackslash{}text\{ N\}\textbackslash{}) as shown in figure. The coefficient of kinetic friction between the blade and the stone is \textbackslash{}(0.60\textbackslash{}), and there is a constant friction torque of \textbackslash{}(6.50\textbackslash{}text\{ N\}\textbackslash{}cdot\textbackslash{}text\{m\}\textbackslash{}) between the axle of the stone and its bearings.

\#\# Question

How much time does it take the grindstone to come from \textbackslash{}(120\textbackslash{}text\{ rev/min\}\textbackslash{}) to rest if it is acted on by the axle friction alone?

\#\# Answer choices

- A: A: 4.86s

- B: B: 3.27s

- C: C: 5.23s

- D: D: 6.86s

\#\# Figure caption (auxiliary; use the image as primary evidence)

The image shows a person sharpening a tool using a large grindstone. Here's a detailed description:

- **Grindstone**: 
  - It's depicted as a large, circular wheel mounted on a frame.
  - It has a mass \textbackslash{}( m = 50.0 \textbackslash{}, \textbackslash{}text\{kg\} \textbackslash{}).

- **Tool and Force**:
  - A hand is applying force \textbackslash{}( F = 160 \textbackslash{}, \textbackslash{}text\{N\} \textbackslash{}) to the tool against the grindstone.
  - This force is represented with a large red arrow pointing towards the grindstone.

- **Rotation**:
  - The grindstone is rotating, indicated by a curved arrow with a symbol \textbackslash{}( \textbackslash{}omega \textbackslash{}), suggesting angular velocity.

- **Frame**:
  - The grindstone is supported by a sturdy, angled frame, likely designed to hold it in place while in use.

- **Other Elements**:
  - The hands and arms appear to be engaging with the tool and the grindstone, illustrating an active use scenario.

- **Overall Context**:
  - The image seems to illustrate a physics or engineering concept related to rotational motion or force application.

\#\# Dataset metadata

Rotational Motion; Multi-Formula Reasoning, Physical Model Grounding Reasoning

\paragraph{Recorded answer/context.}
B: B: 3.27s

\paragraph{Figure/image path.}
/root/proposal\_for\_physic/science-mango/phyx\_mini\_run/phyx\_data/test\_image/781.png

\paragraph{Formalization target.}
create a compiling Lean file with sorry bodies at `PhyXMiniProblems/problem\_phyx\_mini\_0781.lean`. The Lean declarations must preserve the physical quantities, dimensions or dimensional roles, figure labels, governing-law hypotheses, and final relation expressed by this problem.
Use Mathlib/Physlib names found through LeanExplore where available. If a physics API is missing, introduce faithful local abstractions rather than scalar placeholder aliases.

\begin{theorem}[Physics formalization target]
\label{thm:physics:phyx_mini_0781:target}
\lean{PhyXMiniProblems.ProblemPhyXMini0781.stoppingDuration_matches_recordedAnswerB}
\uses{def:physics:phyx-mini-0781:phyxminiproblems-problemphyxmini0781-timeinseconds, def:physics:phyx-mini-0781:phyxminiproblems-problemphyxmini0781-grindstonestoppingsetup, def:physics:phyx-mini-0781:phyxminiproblems-problemphyxmini0781-matchesprimarygrindstonefigure, def:physics:phyx-mini-0781:phyxminiproblems-problemphyxmini0781-matchesgrindstoneproblemdata, def:physics:phyx-mini-0781:phyxminiproblems-problemphyxmini0781-hasphysicalgrindstoneparameters, def:physics:phyx-mini-0781:phyxminiproblems-problemphyxmini0781-satisfiesgrindstonerotationallaws, def:physics:phyx-mini-0781:phyxminiproblems-problemphyxmini0781-recordeddatasetanswer, def:physics:phyx-mini-0781:phyxminiproblems-problemphyxmini0781-matchesdisplayedstoppingtime}
The assigned autoformalize agent should translate this physics problem into Lean declarations in the covered file, with theorem and lemma proof bodies written as `by sorry`.
\end{theorem}
\begin{proof}
This is an autoformalization task, not a proof task. Produce statements that can later be proved by the physics prover without weakening the physical model.
\end{proof}

% --- Archon declaration topology begin ---
\section{Declaration topology}
\begin{definition}[force Dimension]
  \label{def:physics:phyx-mini-0781:phyxminiproblems-problemphyxmini0781-forcedimension}
  \lean{PhyXMiniProblems.ProblemPhyXMini0781.forceDimension}
  The physical dimension of force, M L T\textasciicircum{}-2
\end{definition}
\begin{proof}
  This is the declared model definition of force Dimension.
\end{proof}
\begin{definition}[torque Dimension]
  \label{def:physics:phyx-mini-0781:phyxminiproblems-problemphyxmini0781-torquedimension}
  \lean{PhyXMiniProblems.ProblemPhyXMini0781.torqueDimension}
  The physical dimension of torque and energy, M L\textasciicircum{}2 T\textasciicircum{}-2
\end{definition}
\begin{proof}
  This is the declared model definition of torque Dimension.
\end{proof}
\begin{definition}[moment Of Inertia Dimension]
  \label{def:physics:phyx-mini-0781:phyxminiproblems-problemphyxmini0781-momentofinertiadimension}
  \lean{PhyXMiniProblems.ProblemPhyXMini0781.momentOfInertiaDimension}
  The physical dimension of axial moment of inertia, M L\textasciicircum{}2
\end{definition}
\begin{proof}
  This is the declared model definition of moment Of Inertia Dimension.
\end{proof}
\begin{definition}[Mass Quantity]
  \label{def:physics:phyx-mini-0781:phyxminiproblems-problemphyxmini0781-massquantity}
  \lean{PhyXMiniProblems.ProblemPhyXMini0781.MassQuantity}
  A nonnegative physical mass
\end{definition}
\begin{proof}
  This is the declared model definition of Mass Quantity.
\end{proof}
\begin{definition}[Length Quantity]
  \label{def:physics:phyx-mini-0781:phyxminiproblems-problemphyxmini0781-lengthquantity}
  \lean{PhyXMiniProblems.ProblemPhyXMini0781.LengthQuantity}
  A nonnegative physical length
\end{definition}
\begin{proof}
  This is the declared model definition of Length Quantity.
\end{proof}
\begin{definition}[Force Magnitude Quantity]
  \label{def:physics:phyx-mini-0781:phyxminiproblems-problemphyxmini0781-forcemagnitudequantity}
  \lean{PhyXMiniProblems.ProblemPhyXMini0781.ForceMagnitudeQuantity}
  \uses{def:physics:phyx-mini-0781:phyxminiproblems-problemphyxmini0781-forcedimension}
  A nonnegative physical force magnitude
\end{definition}
\begin{proof}
  This is the declared model definition of Force Magnitude Quantity.
\end{proof}
\begin{definition}[Torque Magnitude Quantity]
  \label{def:physics:phyx-mini-0781:phyxminiproblems-problemphyxmini0781-torquemagnitudequantity}
  \lean{PhyXMiniProblems.ProblemPhyXMini0781.TorqueMagnitudeQuantity}
  \uses{def:physics:phyx-mini-0781:phyxminiproblems-problemphyxmini0781-torquedimension}
  A nonnegative physical torque magnitude
\end{definition}
\begin{proof}
  This is the declared model definition of Torque Magnitude Quantity.
\end{proof}
\begin{definition}[Axial Torque Quantity]
  \label{def:physics:phyx-mini-0781:phyxminiproblems-problemphyxmini0781-axialtorquequantity}
  \lean{PhyXMiniProblems.ProblemPhyXMini0781.AxialTorqueQuantity}
  \uses{def:physics:phyx-mini-0781:phyxminiproblems-problemphyxmini0781-torquedimension}
  A signed axial torque. Positive sign is chosen to agree with the initial rotation arrow; axle friction therefore has negative signed readout
\end{definition}
\begin{proof}
  This is the declared model definition of Axial Torque Quantity.
\end{proof}
\begin{definition}[Time Quantity]
  \label{def:physics:phyx-mini-0781:phyxminiproblems-problemphyxmini0781-timequantity}
  \lean{PhyXMiniProblems.ProblemPhyXMini0781.TimeQuantity}
  A nonnegative physical time interval
\end{definition}
\begin{proof}
  This is the declared model definition of Time Quantity.
\end{proof}
\begin{definition}[Axial Moment Of Inertia Quantity]
  \label{def:physics:phyx-mini-0781:phyxminiproblems-problemphyxmini0781-axialmomentofinertiaquantity}
  \lean{PhyXMiniProblems.ProblemPhyXMini0781.AxialMomentOfInertiaQuantity}
  \uses{def:physics:phyx-mini-0781:phyxminiproblems-problemphyxmini0781-momentofinertiadimension}
  A nonnegative axial moment of inertia
\end{definition}
\begin{proof}
  This is the declared model definition of Axial Moment Of Inertia Quantity.
\end{proof}
\begin{definition}[Axial Angular Velocity Quantity]
  \label{def:physics:phyx-mini-0781:phyxminiproblems-problemphyxmini0781-axialangularvelocityquantity}
  \lean{PhyXMiniProblems.ProblemPhyXMini0781.AxialAngularVelocityQuantity}
  A signed axial angular velocity; radians are dimensionless
\end{definition}
\begin{proof}
  This is the declared model definition of Axial Angular Velocity Quantity.
\end{proof}
\begin{definition}[Axial Angular Acceleration Quantity]
  \label{def:physics:phyx-mini-0781:phyxminiproblems-problemphyxmini0781-axialangularaccelerationquantity}
  \lean{PhyXMiniProblems.ProblemPhyXMini0781.AxialAngularAccelerationQuantity}
  A signed axial angular acceleration
\end{definition}
\begin{proof}
  This is the declared model definition of Axial Angular Acceleration Quantity.
\end{proof}
\begin{definition}[nonnegative SIReadout]
  \label{def:physics:phyx-mini-0781:phyxminiproblems-problemphyxmini0781-nonnegativesireadout}
  \lean{PhyXMiniProblems.ProblemPhyXMini0781.nonnegativeSIReadout}
  Read a nonnegative dimensionful quantity in coherent SI units
\end{definition}
\begin{proof}
  This is the declared model definition of nonnegative SIReadout.
\end{proof}
\begin{definition}[signed SIReadout]
  \label{def:physics:phyx-mini-0781:phyxminiproblems-problemphyxmini0781-signedsireadout}
  \lean{PhyXMiniProblems.ProblemPhyXMini0781.signedSIReadout}
  Read a signed dimensionful quantity in coherent SI units
\end{definition}
\begin{proof}
  This is the declared model definition of signed SIReadout.
\end{proof}
\begin{definition}[mass In Kilograms]
  \label{def:physics:phyx-mini-0781:phyxminiproblems-problemphyxmini0781-massinkilograms}
  \lean{PhyXMiniProblems.ProblemPhyXMini0781.massInKilograms}
  \uses{def:physics:phyx-mini-0781:phyxminiproblems-problemphyxmini0781-massquantity, def:physics:phyx-mini-0781:phyxminiproblems-problemphyxmini0781-nonnegativesireadout}
  Kilogram readout of a physical mass
\end{definition}
\begin{proof}
  This is the declared model definition of mass In Kilograms.
\end{proof}
\begin{definition}[length In Meters]
  \label{def:physics:phyx-mini-0781:phyxminiproblems-problemphyxmini0781-lengthinmeters}
  \lean{PhyXMiniProblems.ProblemPhyXMini0781.lengthInMeters}
  \uses{def:physics:phyx-mini-0781:phyxminiproblems-problemphyxmini0781-lengthquantity, def:physics:phyx-mini-0781:phyxminiproblems-problemphyxmini0781-nonnegativesireadout}
  Meter readout of a physical length
\end{definition}
\begin{proof}
  This is the declared model definition of length In Meters.
\end{proof}
\begin{definition}[force In Newtons]
  \label{def:physics:phyx-mini-0781:phyxminiproblems-problemphyxmini0781-forceinnewtons}
  \lean{PhyXMiniProblems.ProblemPhyXMini0781.forceInNewtons}
  \uses{def:physics:phyx-mini-0781:phyxminiproblems-problemphyxmini0781-forcemagnitudequantity, def:physics:phyx-mini-0781:phyxminiproblems-problemphyxmini0781-nonnegativesireadout}
  Newton readout of a force magnitude
\end{definition}
\begin{proof}
  This is the declared model definition of force In Newtons.
\end{proof}
\begin{definition}[torque Magnitude In Newton Meters]
  \label{def:physics:phyx-mini-0781:phyxminiproblems-problemphyxmini0781-torquemagnitudeinnewtonmeters}
  \lean{PhyXMiniProblems.ProblemPhyXMini0781.torqueMagnitudeInNewtonMeters}
  \uses{def:physics:phyx-mini-0781:phyxminiproblems-problemphyxmini0781-torquemagnitudequantity, def:physics:phyx-mini-0781:phyxminiproblems-problemphyxmini0781-nonnegativesireadout}
  Newton-meter readout of a torque magnitude
\end{definition}
\begin{proof}
  This is the declared model definition of torque Magnitude In Newton Meters.
\end{proof}
\begin{definition}[axial Torque In Newton Meters]
  \label{def:physics:phyx-mini-0781:phyxminiproblems-problemphyxmini0781-axialtorqueinnewtonmeters}
  \lean{PhyXMiniProblems.ProblemPhyXMini0781.axialTorqueInNewtonMeters}
  \uses{def:physics:phyx-mini-0781:phyxminiproblems-problemphyxmini0781-axialtorquequantity, def:physics:phyx-mini-0781:phyxminiproblems-problemphyxmini0781-signedsireadout}
  Signed newton-meter readout of an axial torque
\end{definition}
\begin{proof}
  This is the declared model definition of axial Torque In Newton Meters.
\end{proof}
\begin{definition}[time In Seconds]
  \label{def:physics:phyx-mini-0781:phyxminiproblems-problemphyxmini0781-timeinseconds}
  \lean{PhyXMiniProblems.ProblemPhyXMini0781.timeInSeconds}
  \uses{def:physics:phyx-mini-0781:phyxminiproblems-problemphyxmini0781-timequantity, def:physics:phyx-mini-0781:phyxminiproblems-problemphyxmini0781-nonnegativesireadout}
  Second readout of a physical duration
\end{definition}
\begin{proof}
  This is the declared model definition of time In Seconds.
\end{proof}
\begin{definition}[moment Of Inertia In Kilogram Meters Squared]
  \label{def:physics:phyx-mini-0781:phyxminiproblems-problemphyxmini0781-momentofinertiainkilogrammeterssquared}
  \lean{PhyXMiniProblems.ProblemPhyXMini0781.momentOfInertiaInKilogramMetersSquared}
  \uses{def:physics:phyx-mini-0781:phyxminiproblems-problemphyxmini0781-axialmomentofinertiaquantity, def:physics:phyx-mini-0781:phyxminiproblems-problemphyxmini0781-nonnegativesireadout}
  Kilogram-meter-squared readout of an axial moment of inertia
\end{definition}
\begin{proof}
  This is the declared model definition of moment Of Inertia In Kilogram Meters Squared.
\end{proof}
\begin{definition}[angular Velocity In Radians Per Second]
  \label{def:physics:phyx-mini-0781:phyxminiproblems-problemphyxmini0781-angularvelocityinradianspersecond}
  \lean{PhyXMiniProblems.ProblemPhyXMini0781.angularVelocityInRadiansPerSecond}
  \uses{def:physics:phyx-mini-0781:phyxminiproblems-problemphyxmini0781-axialangularvelocityquantity, def:physics:phyx-mini-0781:phyxminiproblems-problemphyxmini0781-signedsireadout}
  Radian-per-second readout of a signed angular velocity
\end{definition}
\begin{proof}
  This is the declared model definition of angular Velocity In Radians Per Second.
\end{proof}
\begin{definition}[angular Velocity In Revolutions Per Minute]
  \label{def:physics:phyx-mini-0781:phyxminiproblems-problemphyxmini0781-angularvelocityinrevolutionsperminute}
  \lean{PhyXMiniProblems.ProblemPhyXMini0781.angularVelocityInRevolutionsPerMinute}
  \uses{def:physics:phyx-mini-0781:phyxminiproblems-problemphyxmini0781-axialangularvelocityquantity, def:physics:phyx-mini-0781:phyxminiproblems-problemphyxmini0781-angularvelocityinradianspersecond}
  Revolutions-per-minute readout. One revolution is 2 * pi radians and one minute is 60 seconds, so the conversion factor is 30 / pi
\end{definition}
\begin{proof}
  This is the declared model definition of angular Velocity In Revolutions Per Minute.
\end{proof}
\begin{definition}[angular Acceleration In Radians Per Second Squared]
  \label{def:physics:phyx-mini-0781:phyxminiproblems-problemphyxmini0781-angularaccelerationinradianspersecondsquared}
  \lean{PhyXMiniProblems.ProblemPhyXMini0781.angularAccelerationInRadiansPerSecondSquared}
  \uses{def:physics:phyx-mini-0781:phyxminiproblems-problemphyxmini0781-axialangularaccelerationquantity, def:physics:phyx-mini-0781:phyxminiproblems-problemphyxmini0781-signedsireadout}
  Radian-per-second-squared readout of a signed angular acceleration
\end{definition}
\begin{proof}
  This is the declared model definition of angular Acceleration In Radians Per Second Squared.
\end{proof}
\begin{definition}[kinetic Friction Coefficient Readout]
  \label{def:physics:phyx-mini-0781:phyxminiproblems-problemphyxmini0781-kineticfrictioncoefficientreadout}
  \lean{PhyXMiniProblems.ProblemPhyXMini0781.kineticFrictionCoefficientReadout}
  Real readout of the dimensionless kinetic-friction coefficient
\end{definition}
\begin{proof}
  This is the declared model definition of kinetic Friction Coefficient Readout.
\end{proof}
\begin{definition}[Grindstone Shape]
  \label{def:physics:phyx-mini-0781:phyxminiproblems-problemphyxmini0781-grindstoneshape}
  \lean{PhyXMiniProblems.ProblemPhyXMini0781.GrindstoneShape}
  The rigid-body idealization stated for the grindstone
\end{definition}
\begin{proof}
  This is the declared model definition of Grindstone Shape.
\end{proof}
\begin{definition}[Grindstone Figure Component]
  \label{def:physics:phyx-mini-0781:phyxminiproblems-problemphyxmini0781-grindstonefigurecomponent}
  \lean{PhyXMiniProblems.ProblemPhyXMini0781.GrindstoneFigureComponent}
  Components visible in primary image 781.png
\end{definition}
\begin{proof}
  This is the declared model definition of Grindstone Figure Component.
\end{proof}
\begin{definition}[Grindstone Figure Label]
  \label{def:physics:phyx-mini-0781:phyxminiproblems-problemphyxmini0781-grindstonefigurelabel}
  \lean{PhyXMiniProblems.ProblemPhyXMini0781.GrindstoneFigureLabel}
  Text and symbol labels printed in primary image 781.png
\end{definition}
\begin{proof}
  This is the declared model definition of Grindstone Figure Label.
\end{proof}
\begin{definition}[Blade Contact Location]
  \label{def:physics:phyx-mini-0781:phyxminiproblems-problemphyxmini0781-bladecontactlocation}
  \lean{PhyXMiniProblems.ProblemPhyXMini0781.BladeContactLocation}
  The contact location shown for the blade and grindstone
\end{definition}
\begin{proof}
  This is the declared model definition of Blade Contact Location.
\end{proof}
\begin{definition}[Applied Force Arrow Direction]
  \label{def:physics:phyx-mini-0781:phyxminiproblems-problemphyxmini0781-appliedforcearrowdirection}
  \lean{PhyXMiniProblems.ProblemPhyXMini0781.AppliedForceArrowDirection}
  Drawing direction of the red applied-force arrow
\end{definition}
\begin{proof}
  This is the declared model definition of Applied Force Arrow Direction.
\end{proof}
\begin{definition}[Rotation Arrow Top Direction]
  \label{def:physics:phyx-mini-0781:phyxminiproblems-problemphyxmini0781-rotationarrowtopdirection}
  \lean{PhyXMiniProblems.ProblemPhyXMini0781.RotationArrowTopDirection}
  Drawing direction of the curved arrow at the top of the disk. This records the primary pixels without assigning an ambiguous front/back viewing convention to the physical axle orientation
\end{definition}
\begin{proof}
  This is the declared model definition of Rotation Arrow Top Direction.
\end{proof}
\begin{definition}[Grindstone Figure]
  \label{def:physics:phyx-mini-0781:phyxminiproblems-problemphyxmini0781-grindstonefigure}
  \lean{PhyXMiniProblems.ProblemPhyXMini0781.GrindstoneFigure}
  \uses{def:physics:phyx-mini-0781:phyxminiproblems-problemphyxmini0781-grindstonefigurecomponent, def:physics:phyx-mini-0781:phyxminiproblems-problemphyxmini0781-grindstonefigurelabel, def:physics:phyx-mini-0781:phyxminiproblems-problemphyxmini0781-bladecontactlocation, def:physics:phyx-mini-0781:phyxminiproblems-problemphyxmini0781-appliedforcearrowdirection, def:physics:phyx-mini-0781:phyxminiproblems-problemphyxmini0781-rotationarrowtopdirection}
  Qualitative and label evidence transcribed from primary image 781.png
\end{definition}
\begin{proof}
  This is the declared model definition of Grindstone Figure.
\end{proof}
\begin{definition}[Grindstone Stopping Setup]
  \label{def:physics:phyx-mini-0781:phyxminiproblems-problemphyxmini0781-grindstonestoppingsetup}
  \lean{PhyXMiniProblems.ProblemPhyXMini0781.GrindstoneStoppingSetup}
  \uses{def:physics:phyx-mini-0781:phyxminiproblems-problemphyxmini0781-massquantity, def:physics:phyx-mini-0781:phyxminiproblems-problemphyxmini0781-lengthquantity, def:physics:phyx-mini-0781:phyxminiproblems-problemphyxmini0781-forcemagnitudequantity, def:physics:phyx-mini-0781:phyxminiproblems-problemphyxmini0781-torquemagnitudequantity, def:physics:phyx-mini-0781:phyxminiproblems-problemphyxmini0781-axialtorquequantity, def:physics:phyx-mini-0781:phyxminiproblems-problemphyxmini0781-timequantity, def:physics:phyx-mini-0781:phyxminiproblems-problemphyxmini0781-axialmomentofinertiaquantity, def:physics:phyx-mini-0781:phyxminiproblems-problemphyxmini0781-axialangularvelocityquantity, def:physics:phyx-mini-0781:phyxminiproblems-problemphyxmini0781-axialangularaccelerationquantity, def:physics:phyx-mini-0781:phyxminiproblems-problemphyxmini0781-grindstoneshape, def:physics:phyx-mini-0781:phyxminiproblems-problemphyxmini0781-grindstonefigure}
  Physical quantities for the full illustrated scenario and for the axle-only stopping interval. The stopping duration and final angular velocity are independent observables: neither is defined from a displayed answer
\end{definition}
\begin{proof}
  This is the declared model definition of Grindstone Stopping Setup.
\end{proof}
\begin{definition}[Matches Primary Grindstone Figure]
  \label{def:physics:phyx-mini-0781:phyxminiproblems-problemphyxmini0781-matchesprimarygrindstonefigure}
  \lean{PhyXMiniProblems.ProblemPhyXMini0781.MatchesPrimaryGrindstoneFigure}
  \uses{def:physics:phyx-mini-0781:phyxminiproblems-problemphyxmini0781-grindstonestoppingsetup}
  Facts visible in the primary image. The image contains no diameter, friction coefficient, axle-torque, initial-rpm, or stopping-time label; those enter separately as verbal problem data
\end{definition}
\begin{proof}
  This is the declared model definition of Matches Primary Grindstone Figure.
\end{proof}
\begin{definition}[Matches Grindstone Problem Data]
  \label{def:physics:phyx-mini-0781:phyxminiproblems-problemphyxmini0781-matchesgrindstoneproblemdata}
  \lean{PhyXMiniProblems.ProblemPhyXMini0781.MatchesGrindstoneProblemData}
  \uses{def:physics:phyx-mini-0781:phyxminiproblems-problemphyxmini0781-massinkilograms, def:physics:phyx-mini-0781:phyxminiproblems-problemphyxmini0781-lengthinmeters, def:physics:phyx-mini-0781:phyxminiproblems-problemphyxmini0781-forceinnewtons, def:physics:phyx-mini-0781:phyxminiproblems-problemphyxmini0781-torquemagnitudeinnewtonmeters, def:physics:phyx-mini-0781:phyxminiproblems-problemphyxmini0781-angularvelocityinradianspersecond, def:physics:phyx-mini-0781:phyxminiproblems-problemphyxmini0781-angularvelocityinrevolutionsperminute, def:physics:phyx-mini-0781:phyxminiproblems-problemphyxmini0781-kineticfrictioncoefficientreadout, def:physics:phyx-mini-0781:phyxminiproblems-problemphyxmini0781-grindstonestoppingsetup}
  Numerical and categorical data supplied by the prose. The blade quantities belong to the full illustrated scenario, while the asked stopping interval explicitly removes blade contact and retains axle friction alone
\end{definition}
\begin{proof}
  This is the declared model definition of Matches Grindstone Problem Data.
\end{proof}
\begin{definition}[Has Physical Grindstone Parameters]
  \label{def:physics:phyx-mini-0781:phyxminiproblems-problemphyxmini0781-hasphysicalgrindstoneparameters}
  \lean{PhyXMiniProblems.ProblemPhyXMini0781.HasPhysicalGrindstoneParameters}
  \uses{def:physics:phyx-mini-0781:phyxminiproblems-problemphyxmini0781-massinkilograms, def:physics:phyx-mini-0781:phyxminiproblems-problemphyxmini0781-lengthinmeters, def:physics:phyx-mini-0781:phyxminiproblems-problemphyxmini0781-forceinnewtons, def:physics:phyx-mini-0781:phyxminiproblems-problemphyxmini0781-torquemagnitudeinnewtonmeters, def:physics:phyx-mini-0781:phyxminiproblems-problemphyxmini0781-timeinseconds, def:physics:phyx-mini-0781:phyxminiproblems-problemphyxmini0781-momentofinertiainkilogrammeterssquared, def:physics:phyx-mini-0781:phyxminiproblems-problemphyxmini0781-angularvelocityinradianspersecond, def:physics:phyx-mini-0781:phyxminiproblems-problemphyxmini0781-angularaccelerationinradianspersecondsquared, def:physics:phyx-mini-0781:phyxminiproblems-problemphyxmini0781-kineticfrictioncoefficientreadout, def:physics:phyx-mini-0781:phyxminiproblems-problemphyxmini0781-grindstonestoppingsetup}
  Positivity and sign conditions for the physical branch in which the stone initially rotates in the chosen positive direction and axle friction slows it. No numerical stopping duration occurs here
\end{definition}
\begin{proof}
  This is the declared model definition of Has Physical Grindstone Parameters.
\end{proof}
\begin{definition}[Satisfies Grindstone Rotational Laws]
  \label{def:physics:phyx-mini-0781:phyxminiproblems-problemphyxmini0781-satisfiesgrindstonerotationallaws}
  \lean{PhyXMiniProblems.ProblemPhyXMini0781.SatisfiesGrindstoneRotationalLaws}
  \uses{def:physics:phyx-mini-0781:phyxminiproblems-problemphyxmini0781-massinkilograms, def:physics:phyx-mini-0781:phyxminiproblems-problemphyxmini0781-lengthinmeters, def:physics:phyx-mini-0781:phyxminiproblems-problemphyxmini0781-forceinnewtons, def:physics:phyx-mini-0781:phyxminiproblems-problemphyxmini0781-torquemagnitudeinnewtonmeters, def:physics:phyx-mini-0781:phyxminiproblems-problemphyxmini0781-axialtorqueinnewtonmeters, def:physics:phyx-mini-0781:phyxminiproblems-problemphyxmini0781-timeinseconds, def:physics:phyx-mini-0781:phyxminiproblems-problemphyxmini0781-momentofinertiainkilogrammeterssquared, def:physics:phyx-mini-0781:phyxminiproblems-problemphyxmini0781-angularvelocityinradianspersecond, def:physics:phyx-mini-0781:phyxminiproblems-problemphyxmini0781-angularaccelerationinradianspersecondsquared, def:physics:phyx-mini-0781:phyxminiproblems-problemphyxmini0781-kineticfrictioncoefficientreadout, def:physics:phyx-mini-0781:phyxminiproblems-problemphyxmini0781-grindstonestoppingsetup}
  The radius and uniform-disk inertia relations, the two unused reference blade friction laws, and the axle-only rotational dynamics. The last relation is the constant-angular-acceleration kinematic law from the independent initial and final angular velocities. No field states the requested stopping time or mentions an answer choice
\end{definition}
\begin{proof}
  This is the declared model definition of Satisfies Grindstone Rotational Laws.
\end{proof}
\begin{lemma}[moment Of Inertia About Axle is 169 over 100]
  \label{lem:physics:phyx-mini-0781:phyxminiproblems-problemphyxmini0781-momentofinertiaaboutaxle-is-169-over-100}
  \lean{PhyXMiniProblems.ProblemPhyXMini0781.momentOfInertiaAboutAxle_is_169_over_100}
  \uses{def:physics:phyx-mini-0781:phyxminiproblems-problemphyxmini0781-momentofinertiainkilogrammeterssquared, def:physics:phyx-mini-0781:phyxminiproblems-problemphyxmini0781-grindstonestoppingsetup, def:physics:phyx-mini-0781:phyxminiproblems-problemphyxmini0781-matchesgrindstoneproblemdata, def:physics:phyx-mini-0781:phyxminiproblems-problemphyxmini0781-satisfiesgrindstonerotationallaws}
  The stated uniform disk has axial moment of inertia 1.69 kg m\textasciicircum{}2
\end{lemma}
\begin{proof}
  The conclusion follows from the governing relations and figure readouts named in the statement of moment Of Inertia About Axle is 169 over 100.
\end{proof}
\begin{lemma}[initial Angular Velocity is four pi]
  \label{lem:physics:phyx-mini-0781:phyxminiproblems-problemphyxmini0781-initialangularvelocity-is-four-pi}
  \lean{PhyXMiniProblems.ProblemPhyXMini0781.initialAngularVelocity_is_four_pi}
  \uses{def:physics:phyx-mini-0781:phyxminiproblems-problemphyxmini0781-angularvelocityinradianspersecond, def:physics:phyx-mini-0781:phyxminiproblems-problemphyxmini0781-grindstonestoppingsetup, def:physics:phyx-mini-0781:phyxminiproblems-problemphyxmini0781-matchesgrindstoneproblemdata}
  The initial 120 rev/min equals 4 * pi rad/s
\end{lemma}
\begin{proof}
  The conclusion follows from the governing relations and figure readouts named in the statement of initial Angular Velocity is four pi.
\end{proof}
\begin{lemma}[axle Only Angular Acceleration is neg 650 over 169]
  \label{lem:physics:phyx-mini-0781:phyxminiproblems-problemphyxmini0781-axleonlyangularacceleration-is-neg-650-over-169}
  \lean{PhyXMiniProblems.ProblemPhyXMini0781.axleOnlyAngularAcceleration_is_neg_650_over_169}
  \uses{def:physics:phyx-mini-0781:phyxminiproblems-problemphyxmini0781-angularaccelerationinradianspersecondsquared, def:physics:phyx-mini-0781:phyxminiproblems-problemphyxmini0781-grindstonestoppingsetup, def:physics:phyx-mini-0781:phyxminiproblems-problemphyxmini0781-matchesgrindstoneproblemdata, def:physics:phyx-mini-0781:phyxminiproblems-problemphyxmini0781-hasphysicalgrindstoneparameters, def:physics:phyx-mini-0781:phyxminiproblems-problemphyxmini0781-satisfiesgrindstonerotationallaws}
  The axle-only angular acceleration is the negative torque magnitude divided by the solid-disk inertia, namely -650 / 169 rad/s\textasciicircum{}2
\end{lemma}
\begin{proof}
  The conclusion follows from the governing relations and figure readouts named in the statement of axle Only Angular Acceleration is neg 650 over 169.
\end{proof}
\begin{definition}[Answer Choice]
  \label{def:physics:phyx-mini-0781:phyxminiproblems-problemphyxmini0781-answerchoice}
  \lean{PhyXMiniProblems.ProblemPhyXMini0781.AnswerChoice}
  Labels of the four stopping-time choices printed with the problem
\end{definition}
\begin{proof}
  This is the declared model definition of Answer Choice.
\end{proof}
\begin{definition}[seconds]
  \label{def:physics:phyx-mini-0781:phyxminiproblems-problemphyxmini0781-answerchoice-seconds}
  \lean{PhyXMiniProblems.ProblemPhyXMini0781.AnswerChoice.seconds}
  \uses{def:physics:phyx-mini-0781:phyxminiproblems-problemphyxmini0781-answerchoice}
  Time in seconds printed beside each answer label
\end{definition}
\begin{proof}
  This is the declared model definition of seconds.
\end{proof}
\begin{definition}[recorded Dataset Answer]
  \label{def:physics:phyx-mini-0781:phyxminiproblems-problemphyxmini0781-recordeddatasetanswer}
  \lean{PhyXMiniProblems.ProblemPhyXMini0781.recordedDatasetAnswer}
  \uses{def:physics:phyx-mini-0781:phyxminiproblems-problemphyxmini0781-answerchoice}
  The answer label recorded in the supplied dataset metadata
\end{definition}
\begin{proof}
  This is the declared model definition of recorded Dataset Answer.
\end{proof}
\begin{definition}[Matches Displayed Stopping Time]
  \label{def:physics:phyx-mini-0781:phyxminiproblems-problemphyxmini0781-matchesdisplayedstoppingtime}
  \lean{PhyXMiniProblems.ProblemPhyXMini0781.MatchesDisplayedStoppingTime}
  \uses{def:physics:phyx-mini-0781:phyxminiproblems-problemphyxmini0781-timequantity, def:physics:phyx-mini-0781:phyxminiproblems-problemphyxmini0781-timeinseconds, def:physics:phyx-mini-0781:phyxminiproblems-problemphyxmini0781-answerchoice}
  A duration agrees with a displayed value to the nearest hundredth of a second when the absolute error is strictly less than half a hundredth
\end{definition}
\begin{proof}
  This is the declared model definition of Matches Displayed Stopping Time.
\end{proof}
% --- Archon declaration topology end ---
% --- Archon physics formalization source end ---
